Use the same ring and scale for the ciphertext and the integer lift of $\Lambda$. Assume $t\mid q$, $\Delta=q/t$, and $\|E\|_\infty<\Delta/2$ for the stated decoded plaintext modulo $t$.

Suppose we have a GLWE ciphertext \textsf{ct} and a new plaintext polynomial $\Lambda$ as follows: 


$\textsf{ct} = \textsf{GLWE}_{S, \sigma}(\Delta M + E) = ( A_0, A_1, \ldots, A_{k-1}, B) \in \mathcal{R}_{\langle n,q \rangle}^{k + 1}$

$\Lambda$: a new plaintext polynomial

$\Delta \Lambda$: a $\Delta$-scaled new plaintext polynomial

$ $

\noindent Let's define the following ciphertext-to-plaintext addition operation: 

$\textsf{ct} + \Delta\Lambda = ( A_0, \text{ } A_1, \ldots, A_{k-1}, \text{ } B + \Delta \Lambda)$

$ $

\noindent Then, the following is true:

\begin{tcolorbox}[title={\textbf{\tboxlabel{\ref*{sec:glwe-add-plain}} GLWE Homomorphic Addition with a Plaintext}}]
$\textsf{GLWE}_{S, \sigma}(\Delta M + E) + \Delta\Lambda $

$=  (\{A_i\}_{i=0}^{k-1}, \text{ } B) + \Delta\Lambda$

$=  (\{A_i\}_{i=0}^{k-1}, \text{ } B + \Delta\Lambda)$

$= \textsf{GLWE}_{S, \sigma}(\Delta (M + \Lambda) + E)$
\end{tcolorbox}


This means that adding a ($\Delta$-scaled) plaintext polynomial $\Lambda$ to a GLWE ciphertext that encrypts $M$ and decrypting it yields $M + \Lambda$. 




\begin{myproof}

\begin{enumerate}
\item Since $B = \sum\limits_{i=0}^{k-1}{(A_i \cdot S_i)} + \Delta \cdot M + E$,

$B + \Delta \cdot \Lambda = \sum\limits_{i=0}^{k-1}{(A_i \cdot S_i)} + \Delta \cdot M + E + \Delta \cdot \Lambda = \sum\limits_{i=0}^{k-1}{(A_i \cdot S_i)} + \Delta \cdot (M + \Lambda) + E$  \\

This means that $(A_0, A_1, ... \text{ } A_{k-1}
, B + \Delta \Lambda)$ form the ciphertext $\textsf{GLWE}_{S, \sigma}(\Delta (\Lambda + M) + E)$
\item Thus, \\
$\textsf{GLWE}_{S, \sigma}(\Delta M + E) + \Delta\Lambda$ \\
$ = (A_0, \text { } A_1, ... \text{ } A_{k-1}, \text { } B + \Delta \Lambda)$ \\
$ = (\{A_i\}_{i=0}^{k-1}, \text { } B + \Delta \Lambda)$ \\
$= \textsf{GLWE}_{S, \sigma}(\Delta (M + \Lambda) + E)$

%\begin{flushright}
%\qedsymbol{} 
%\end{flushright}
\end{enumerate}

\end{myproof}


\para{Noise Growth:} The phase error relative to the unreduced message $M+\Lambda$ remains $E$. With $t\mid q$, reducing that message modulo $t$ does not alter the effective error. If $t\nmid q$, a message wrap-around can add the overflow-dependent error described in \autoref{subsubsec:lwe-noise-bound}; the statement that noise stays unchanged must then be understood relative to the unreduced message.